

\illustration{0.30}{circuits}{\href{https://unsplash.com/photos/jXd2FSvcRr8}{Picture} by Umberto, Unsplash Licence}%

Sara is doing her summer internship at NCPC (Never Crashing Personal Computers). One day, a rare creature
appeared in the office: an algorithmic problem!

The company has a machine that places electronic components on circuit boards. Normally, it would 
do this one component at a time. But recently the machine has received an update which allows it 
to place two different components simultaneously. The bottleneck then becomes the component with greater placement time. 
Now it is far from obvious what strategy the machine should use in order to minimize the total placement
time. Sara decides to write an algorithm to determine this strategy. 

You have $N$ different types of electronic components. There are $f_i$ copies of the $i$th type, and 
the components of this type have a placement time of $t_i$ nanoseconds.
The goal is to place all of the components using a sequence of moves.
In one move, the machine can take two components of type $i$ and $j$, where $i \neq j$, and place 
both of them simultaneously. This takes $\max(t_i, t_j)$ nanoseconds. 
The machine can also place a single component of type $i$ in one move, which takes $t_i$ nanoseconds.

Calculate the minimum possible time to place all components.

